\section{Overview}
\label{m:sec:overview}

The processes studied in this thesis describe populations small enough that their
fate is decided by chance rather than by their mean behaviour --- a founding cohort
of a few individuals, or the contents of a single compartment that may later
rupture. At those counts a deterministic rate equation answers the wrong question.
It returns a single trajectory, growing or decaying, when what one wants to know is
whether the population establishes itself at all, and how large it is if it does.
Both are questions about a distribution, and neither is visible in a mean.

This chapter assembles the mathematical framework in which those questions are
posed, together with the methods that answer them. The two are kept apart on
purpose. Naming the processes first, and only then developing each technique once,
on whichever process displays it most clearly, avoids the usual awkwardness of a
background chapter, where the same method is introduced three times over because
three models happen to need it.

\Cref{m:sec:chains} sets out the \emph{objects}: the classes of Markov process that
appear in this thesis. Discrete-time chains come first, the simple random walk as
the elementary example and the Galton--Watson branching process as the one that
matters here. Continuous-time chains follow, built from exponential clocks: the
Poisson process, the linear birth--death process, and the
birth--death--catastrophe process in which an abrupt event can remove an entire
population at a stroke. Two extensions are then described that leave the
time-homogeneous setting: processes whose rates vary with time, worked through on a
logistic model of diversification in which the birth rate declines as a niche fills;
and coupled systems in which a deterministic differential equation and a Markov chain
each drive the other, whose schematic instance here is a compartment that ruptures
into a shared medium.

\Cref{m:sec:methods} sets out the \emph{methods}. Discrete time is treated briefly,
because the technique that the branching problem of this thesis actually requires
--- the Koenigs linearisation of a functional iteration --- is developed in \ChA{}
rather than here. Continuous time is treated at length, with the birth--death
process as the running example, and covers backward equations, means, variances and
higher moments, hitting probabilities, extinction probabilities, and the limiting
conditional means that arise when one conditions on a process not yet having been
absorbed. \Cref{m:sec:moc} then develops the method of characteristics, which
converts the master equation for an infinite family of state probabilities into a
single first-order partial differential equation for a generating function and
solves it; a worked example is carried through in full.

Three threads recur, and it is worth naming them at the outset.

The first is \emph{conditioning on survival}. A subcritical branching process dies
out with probability one, so its unconditional mean decays geometrically and
records little beyond the certainty of extinction. Restrict attention to those
realisations still alive at generation $n$, however, and the picture changes
entirely: the conditional population settles to a fixed distribution, and the mean
of that distribution is a constant. For the binary Galton--Watson process with
division probability $p<\tfrac12$ that constant is $1/A(p)$, where
$A(p) = \lim_n S_n/(2p)^n$ and $S_n$ is the survival probability to generation $n$.
This chapter defines $A(p)$ and shows how the whole limiting conditional law is
expressed through it, in \cref{m:sec:condmeans}. It does not analyse the constant:
the exact representations, the bounds, the near-critical asymptotics, and the
dynamical obstruction that explains why no elementary formula for it has been found
are the subject of \ChA.

The second is the \emph{near-critical regime}. As $p$ rises towards $\tfrac12$ the
process takes ever longer to die, the quasi-stationary mean diverges, and any
numerical scheme that works by iterating to stationarity becomes impractical. Just
above $\tfrac12$ the size of the founding cohort begins to weigh heavily on whether
the process survives at all. The behaviour on either side of $p=\tfrac12$ is what
makes small founding populations interesting, and it is also what makes them
awkward to compute with. The phenomenology is recorded in \cref{m:app:criticalgw};
the theory of $A(p)$ that eventually tames it belongs to \ChA.

The third is \emph{absorption}. Extinction and rupture are both absorbing states,
and much of this thesis concerns a chain conditioned on not yet having reached one.
The generating-function methods of \cref{m:sec:methods} are what make such
conditioning computable.

Two remarks on how the chapter is written. First, not all of it is background: the
closed form obtained in \cref{m:app:absbirthdeath} for the absorption--birth--death
generating function is a result of this thesis, as is the integral identity of
\cref{m:app:hyp} on which an alternative derivation of it rests. Everything else is
standard, and is presented as such. Second, the argument is in places heuristic ---
continuum approximations to difference equations, mean-field substitutions,
asymptotic matching. Every such step is marked where it occurs, and where a claim
can be checked numerically it has been.
